% !TEX root = ../main.tex
\section{Código fuente del motor difuso}
\label{ap:codigo}

Este apéndice incluye los módulos esenciales del paquete \texttt{fuzzy-core} para permitir verificación independiente de la implementación. Cada listado corresponde literalmente al archivo del repositorio.

\subsection{Funciones de pertenencia (\texttt{membership.ts})}

\begin{lstlisting}[language=Java,caption={Funciones triangular y trapezoidal con clamp en [0,1].}]
export function triangular(x, a, b, c) {
  if (x <= a || x >= c) return 0;
  if (x === b) return 1;
  if (x < b) return clamp01((x - a) / (b - a));
  return clamp01((c - x) / (c - b));
}

export function trapezoidal(x, a, b, c, d) {
  if (x < a || x > d) return 0;
  if (a === b && x <= b) return 1;
  if (c === d && x >= c) return 1;
  if (x >= b && x <= c) return 1;
  if (x > a && x < b) return clamp01((x - a) / (b - a));
  if (x > c && x < d) return clamp01((d - x) / (d - c));
  return 0;
}
\end{lstlisting}

\subsection{Inferencia Mamdani (\texttt{mamdani.ts})}

\begin{lstlisting}[language=Java,caption={Construccion de activaciones, agregacion y centroide.}]
const ruleActivations = system.rules.map((rule) => {
  const antecedentDegrees = rule.antecedents.map((a) => ({
    variable: a.variable,
    term: a.term,
    degree: fuzzification[a.variable][a.term] ?? 0,
  }));
  const alpha = Math.min(...antecedentDegrees.map((d) => d.degree));
  const consequent = output.terms.find((t) => t.id === rule.consequent.term);
  const clippedArea = sampleUniverse(output, resolution)
    .reduce((area, x) => area + Math.min(alpha, evaluateMembership(consequent.shape, x)) * resolution, 0);
  return { rule, alpha, antecedentDegrees, clippedArea };
});

for (let x = output.min; x <= output.max; x += resolution) {
  const sample = { x, aggregated: 0 };
  for (const term of output.terms) {
    const termAlpha = Math.max(0, ...ruleActivations
      .filter((act) => act.rule.consequent.term === term.id)
      .map((act) => act.alpha));
    sample[term.id] = Math.min(termAlpha, evaluateMembership(term.shape, x));
    sample.aggregated = Math.max(sample.aggregated, sample[term.id]);
  }
  outputSamples.push(sample);
}
\end{lstlisting}

\subsection{Defuzzificacion por centroide (\texttt{centroid.ts})}

\begin{lstlisting}[language=Java,caption={Centroide discreto.}]
export function centroid(samples) {
  const numerator = samples.reduce((s, p) => s + p.x * p.aggregated, 0);
  const denominator = samples.reduce((s, p) => s + p.aggregated, 0);
  return {
    value: denominator === 0 ? 0 : numerator / denominator,
    numerator,
    denominator,
  };
}
\end{lstlisting}

La implementación completa, incluyendo definición de variables, reglas, casos de prueba y configuración del monorepo, se encuentra en el repositorio del proyecto bajo el directorio \texttt{packages/fuzzy-core/src}.
